\section*{TD 5 --- Arbres et ensembles}
\addcontentsline{toc}{section}{TD 5 --- Arbres et ensembles}

\subsection*{Exercice 1 --- Impureté de Gini \hfill \normalfont\footnotesize (12 min)}
Un nœud contient 10 exemples : 6 positifs, 4 négatifs.
\begin{enumerate}
  \item Calculer son indice de Gini.
  \item On le coupe en deux enfants : $A$ (5 positifs, 1 négatif) et $B$
        (1 positif, 3 négatifs). Calculer le Gini de chacun.
  \item Calculer le Gini pondéré des enfants, puis le gain.
  \item La coupe est-elle utile ?
\end{enumerate}

\subsection*{Exercice 2 --- Arbre ou forêt ? \hfill \normalfont\footnotesize (8 min)}
Un arbre unique atteint 0,96 en entraînement et 0,71 en test.
\begin{enumerate}
  \item Quel est le problème ?
  \item Comment une forêt aléatoire le corrige-t-elle ? Expliquer en termes de
        variance.
  \item Pourquoi une forêt ne descend-elle jamais à 0 d'erreur d'entraînement ?
\end{enumerate}

\subsection*{Exercice 3 --- Importance et redondance \hfill \normalfont\footnotesize (8 min)}
Deux colonnes sont corrélées à $r=0{,}95$, et le modèle leur attribue une
importance de 0,20 à chacune, contre 0,45 pour une troisième variable.
\begin{enumerate}
  \item L'importance totale des deux colonnes redondantes est-elle faible ?
  \item Peut-on conclure que chacune compte peu ?
  \item Que faire pour y voir clair ?
\end{enumerate}

\subsection*{Exercice 4 --- Bagging ou boosting ? \hfill \normalfont\footnotesize (8 min)}
Associer chaque affirmation au bon ensemble.
\begin{enumerate}
  \item Les arbres sont entraînés en parallèle, puis moyennés.
  \item Chaque arbre corrige les erreurs du précédent.
  \item Réduit surtout la variance.
  \item Réduit surtout le biais, mais peut surapprendre si l'on ajoute trop
        d'itérations.
\end{enumerate}

\subsection*{Exercice 5 --- Interpréter une décision \hfill \normalfont\footnotesize (6 min)}
Un modèle refuse un crédit. On vous demande d'expliquer la décision à un
client.
\begin{enumerate}
  \item Pourquoi l'importance globale des variables ne suffit-elle pas ?
  \item Que apporte une explication locale (SHAP) ?
  \item Quelle précaution prendre en présentant ces valeurs ?
\end{enumerate}


\subsection*{Exercice 6 --- Entropie contre Gini \hfill \normalfont\footnotesize (8 min)}
Un nœud contient 6 positifs et 4 négatifs.
\begin{enumerate}
  \item Calculer son indice de Gini (déjà vu au TD précédent).
  \item Calculer son entropie $-\sum p_k\log_2 p_k$.
  \item Les deux critères classent-ils les mêmes découpages de la même façon ?
        Que change le choix ?
\end{enumerate}

\subsection*{Exercice 7 --- Profondeur et généralisation \hfill \normalfont\footnotesize (7 min)}
Un arbre donne les résultats suivants :

\begin{center}
\small
\begin{tabular}{@{}lccc@{}}
\toprule
\textbf{Profondeur} & 2 & 5 & 12 \\
\midrule
Erreur entraînement & 0,22 & 0,08 & 0,01 \\
Erreur validation & 0,24 & 0,11 & 0,19 \\
\bottomrule
\end{tabular}
\end{center}

\begin{enumerate}
  \item Quelle profondeur choisir ?
  \item Que se passe-t-il au-delà ?
  \item Comment obtient-on cette courbe sans toucher au test ?
\end{enumerate}

\subsection*{Exercice 8 --- Pourquoi moyenner aide \hfill \normalfont\footnotesize (9 min)}
Cinq arbres ont une variance individuelle $\sigma^2=100$ et une corrélation
moyenne $\rho=0{,}2$.
\begin{enumerate}
  \item Rappeler la variance d'un arbre seul.
  \item Calculer la variance de leur moyenne :
        $\rho\sigma^2+(1-\rho)\sigma^2/5$.
  \item Que se passerait-il avec $\rho=0{,}9$ ? Conclure sur l'intérêt de
        décorréler les arbres.
\end{enumerate}

\subsection*{Exercice 9 --- Régler un boosting \hfill \normalfont\footnotesize (7 min)}
On augmente le nombre d'itérations de 100 à 2\,000.
\begin{enumerate}
  \item Que fait l'erreur d'entraînement ?
  \item Et l'erreur de validation ?
  \item Quel rôle joue le pas d'apprentissage dans ce compromis ?
\end{enumerate}

\newpage
